\subsection{A warranty can promise another year of service}
\label{sec:service-warranty}

A returned product and a warranty claim need different histories. The return
example unwinds a sale when a customer exercises a refund right. A warranty
may instead promise to keep a product working, or provide a service after the
product has been delivered. The company must identify which promise it made
before deciding how much of a customer's payment it has earned. The following
cases use the same retained AASB 15 compilation and historical Australian
for-profit scope as the return example.

Suppose a customer buys a product and a separately purchasable service warranty
for a combined 108. Comparable standalone selling prices at contract inception
are 100 for the product and 20 for the warranty. Both are distinct promises;
the product works independently of the additional service. The contract has
neither a discount attributable to just one promise nor variable consideration.
Under paragraphs 73--86, the 12 bundle discount is therefore allocated in
proportion to the standalone prices. The product receives
$108\times100/120=90$ and the service receives $108\times20/120=18$.
An invoice's labels would not by themselves establish those standalone prices
or justify a different discount allocation \citep{aasb15dated}.

On 1 July 2024 the company receives 108 and transfers control of the product,
whose inventory cost is 60. First, receipt debits cash and credits the contract
liability by 108. Delivery then debits that liability and credits product
revenue by 90, while debiting cost of sales and crediting inventory by 60.
Profit is 30 and a liability of 18 remains. Recognising all 108 on delivery
would credit revenue for service the company still owes.

The service warranty runs for twelve monthly service periods to 30 June 2025.
The company stands ready uniformly; the customer receives and consumes the
benefits as service is supplied. The admitted measure gives each completed
service month one twelfth of the allocation and faithfully depicts performance.
These facts support over-time recognition under
35(a), 39--45 and B14--B19; the existence of a twelve-month date range alone
would not suffice. At 31 December 2024, half the service has been provided,
so half of 18, or 9, is earned. Debit the contract liability and credit
service revenue by 9. Actual service costs of 3 are expensed and paid.
Opening inventory and equity were both 60, with no other opening balances.
Year-one revenue is now 99, expenses 63 and pretax profit 36. Cash 105 equals
the remaining liability 9 plus equity 96.

The second half-year releases the remaining liability of 9 and incurs another
3 of paid service costs. Its profit is 6 and closing cash and equity are both
102. Over the complete contract, revenue 108 less costs 66 gives profit 42.
The cash statement follows the same history: operating cash is 105 in the
first period and $-3$ in the second. The independent accrual bridges are
$36+60+9=105$ and $6-9=-3$, using the inventory reduction and changes in the
contract liability. There is no OCI or subsequent owner transaction in this
pretax slice; the complete company calculation still needs tax and any other
obligations.

The contract record must therefore retain the two promises, the original
allocation and its price evidence, delivery identity, coverage, progress
measure, cumulative service revenue and actual costs. Each service update
posts the difference between newly supported cumulative revenue and revenue
already recognised. Replaying an accepted delivery cannot release the product
allocation twice. Unknown progress is not zero performance: paragraphs 44--45
require a separate assessment of reasonable measurement and, where applicable,
cost recovery. This percentage method cannot invent those findings.

The classification also has a substantive boundary. A separately purchasable
warranty is a distinct service under B29. Without that option, B30 sends an
assurance-only warranty to AASB 137 unless it provides additional service.
Law, duration and promised tasks help distinguish the cases under B31; there
is no universal number of months that makes a warranty a service. If assurance
and additional service cannot reasonably be accounted for separately, B32
treats the two warranties together as one performance obligation. Harm or
infringement compensation under B33 is not such a service. The exact Australian
AASB 137 measurement text has not been inspected for this extension, so an
assurance obligation retains a required measurement method and facts. It is
not assigned a zero provision or converted into the service example.

\subsection{Changed promises can affect future units or past progress}
\label{sec:customer-modification}

A contract can change after some performance has already been recognised.
The accounting consequence depends on what remains to be supplied. Consider
an original prepaid order for ten distinct units at 10 each. Four have been
delivered, earning 40. The parties approve an enforceable amendment adding
six units for 48. Comparable standalone prices remain 10 each, and no
appropriate contract-specific adjustment makes the 48 addition reflect
those prices. The addition therefore fails the separate-contract condition
in paragraph 20. Because the remaining units are distinct from those already
delivered, paragraph 21(a) reallocates the consideration for the remaining
performance \citep{aasb15dated}.

The undelivered original units carry 60 of consideration. Adding the new 48
gives 108 for twelve remaining units, whose standalone prices are equal.
Thus their new unit allocation is
\[
 \frac{\text{unrecognised original consideration}
       +\text{additional consideration}}
      {\text{remaining equal-price units}}
 =\frac{60+48}{6+6}=9.
\]
The first four units keep their revenue of 40. Spreading the entire revised
price over all sixteen units would instead give $148/16=9.25$ and rewrite
their revenue as 37. That calculation is wrong for the stated facts: this
modification reallocates remaining performance, not completed deliveries.

Inventory and cash make the consequences concrete. Start with sixteen units
costing 6 each, financed by opening equity of 96; six units were initially
uncommitted. The first prepayment debits cash and credits a contract liability
by 100. The first four deliveries release 40 of revenue and remove inventory
cost 24. Immediately before amendment, cash 100 plus inventory 72 equals the
liability 60 plus equity 112. Receipt of the additional 48 increases cash
and the liability. Approving the changed order itself earns no new revenue
and creates no new stock.

Deliver six more units by the first year end. They release liability and
revenue of 54, with inventory cost 36. Total year-one revenue is 94, cost
of sales 60 and profit 34. Cash 148 and inventory 36 equal liability 54
and equity 130. Delivering the final six in the next year earns another
54 at cost 36, leaving cash and equity 148 with no inventory or contract
liability. That period's profit is 18 but no cash arrives: the customer
already paid. The operating cash bridges are $34+60+54=148$ and
$18+36-54=0$. Lifetime revenue 148 less inventory cost 96 gives profit 52.

If the additional six units instead cost 60 and meet paragraph 20's
standalone-price condition, the addition is a separate contract. The first
four still earn 40; the six remaining original units earn 60 and the new
six earn 60. Both contracts keep their own allocation records. The difference
between these cases is the admitted price and scope assessment, not permission
to choose whichever result produces the preferred margin.

An integrated promise can have the opposite effect on past cumulative
revenue. Suppose a customer controls work in progress as the company creates
it, satisfying paragraph 35(b). The company has one integrated performance
obligation, and its remaining work is not distinct from work already supplied.
Original price 100 and supported output progress of 40\% yielded revenue 40.
An approved enforceable amendment changes that same promise to price 130;
the revised faithful output measure is 30\% of the expanded work. Under
21(b), both the price and progress change the cumulative amount earned:
\[
 \text{revised cumulative revenue}=130\times0.30=39,
 \qquad \text{catch-up}=39-40=-1.
\]
The new work makes the completed work a smaller share of the promise.
Revenue decreases by 1 at amendment even though the total contract price
increases. The old 40 remains attached to its original information view;
the modification does not establish that the earlier estimate was an error.

For a complete prepaid example, start with cash and equity 50. Receive 100
and pay service costs 30 before amendment. Recognition of 40 leaves cash
120, liability 60 and equity 60. At amendment, receive the additional 30
and debit revenue 1 while crediting the liability 1. Cash is now 150,
the liability is 91 and equity is 59: cumulative revenue 39 less costs 30
has earned profit 9. Operating cash since opening is 100, also $9+91$.
Completing the promise next period releases revenue 91 and incurs further
paid service costs of 60. Period profit is 31, operating cash is
$31-91=-60$, and closing cash and equity are 90. Lifetime profit is
$130-90=40$. The output progress percentages were independently supported;
they were not inferred from these cash costs.

Both modifications need the approved rights, effective date, previous and
revised promise, consideration, performance measurements and source versions.
The distinct-goods method freezes the fulfilled deliveries and replaces the
remaining allocation; the integrated method proposes only the difference
between revised cumulative revenue and the amount already recognised.
Economic state and journals change atomically. Different billing schedules
can require a contract-asset adjustment instead of the prepaid liability
entries used here. A request that has not been approved leaves the existing
contract operative under paragraph 18. Unknown enforceability stays
unresolved; approved scope with undetermined price requires paragraph 19's
variable-consideration estimate and constraint, not a price of zero.

Paragraph 21(c) also admits mixed remaining obligations, requiring an explicit
allocation consistent with the two principles above. The examples do not
force such a contract into one category. Price-estimate changes without a
contract modification follow 87--89. A later price change after modification
needs paragraph 90: variable consideration promised before a 21(a)
modification can still be allocated to the earlier obligations under 90(a).
It cannot automatically be spread over the twelve remaining units merely
because their current allocation is 9 each.

\subsection{The same cash can accompany different revenue}
\label{sec:principal-agent-cash}

Suppose an entity receives 120 from a customer and sends 100 to a supplier.
It keeps 20. Those cash movements cannot tell us whether it sold the supplier's
good or arranged a sale by that supplier. Under B34--B38, the decision starts
with the specified good or service and whether the entity controls it before
transfer to the customer. Responsibility for delivery, inventory risk and
price discretion inform that decision; none is a substitute for control, and
momentary legal title does not by itself establish it \citep{aasb15dated}.

In the principal case, the entity controls the identified good before
customer transfer, can direct its use and obtain its benefits, and is
responsible for delivering it. Acquisition on credit debits inventory 100
and credits the supplier payable 100. Customer transfer and cash receipt
debit cash 120 and credit revenue 120; the inventory leaves through cost
of sales 100. Settling the supplier debits the payable and credits cash 100.
From zero opening balances, the result is cash and equity 20, revenue 120
and profit 20.

In the agent case, the supplier controls the good until transfer; the entity
only arranges the sale and earns a fixed commission of 20 when that service
is completed. Assume the entity controls the collected cash and owes the
supplier 100 under its remittance agreement. Funds held outside its balance
sheet in a fiduciary capacity are outside this example. Receipt therefore
debits cash 120, credits the supplier payable 100 and credits commission
revenue 20. Settlement clears the payable against cash 100. The agent has
no inventory or supplier cost of sales. Its final cash, equity and profit
are again 20, but revenue is 20.

Even the intermediate balance sheets agree: before remittance, both have
cash 120, a payable 100 and equity 20. The different control facts explain
the different revenue. The method must retain those facts and classify each
specified promise separately, since the same entity can act as principal
for one promise and agent for another. An unknown control finding leaves
classification unresolved; the receipt cannot decide it. Both cash legs
remain in the ledger. Whether the cash-flow statement presents receipts and
remittances gross or net requires its own AASB 107 conditions, which do not
follow merely from the revenue classification. These examples establish
the cash change and pretax profit contribution while preserving that separate
presentation question.
